\chapter{The command line}
\label{ch:cli}

The program is \cmd{vet}. It has two commands, \cmd{vet apply} and \cmd{vet init}, and
one option for itself, \opt{--version}. \cmd{python -m hdl_vet} is the same program.
The options \code{-h} and \code{--help} of \cmd{vet}, \cmd{vet apply} and \cmd{vet init}
print a short usage text. This chapter describes every option in full.

Chapter~\ref{ch:first} explains the file list, the base and the input tree.

\section{vet apply}

\cmd{vet apply} runs scripts on a design. These are its options.

\begin{description}
\item[\opt{-f}, \opt{--filelist} FILE] The top-level file list. It is required,
  unless you give \opt{--resume} or the project file names one. It replaces the
  file list and the base of the project.
\item[\opt{-o}, \opt{--outdir} DIR] The folder for the output tree. DIR must be a
  new folder or an empty folder, unless you give \opt{--force}. It must not lie inside
  the input tree, and it must not contain the input tree. It must not contain a script
  that runs or the project file.
\item[\opt{-s}, \opt{--script} FILE] A script. You can give it many times. The scripts run
  in the order given, after the scripts of the project file. A file that you name twice is
  an error.
\item[\opt{--top} MODULE] The top module. It replaces the top module of the project.
  Always name it (chapter~\ref{ch:signoff}).
\item[\opt{--base} DIR] The base of the file list (chapter~\ref{ch:first}). It replaces
  the base of the project.
\item[\opt{--dry-run}] Print the diff of the changed files and write no output tree.
  The sign-off runs. An archive that you name with \opt{--save} is still written.
\item[\opt{--force}] Replace an output folder that exists (below).
\item[\opt{--save} FILE] After the scripts end, save the session to a zip archive
  (chapter~\ref{ch:session}).
\item[\opt{--resume} FILE] Continue from a saved session. hdl-vet does not read a
  file list. Do not give \opt{-f}, \opt{--top} or \opt{--base} with it.
\item[\opt{--project} FILE] Use this project file instead of \file{hdl-vet.toml} in
  the current folder (chapter~\ref{ch:project}).
\item[\opt{--no-project}] Ignore the project file. You cannot give it together with
  \opt{--project}.
\end{description}

You must give at least one of \opt{-o}, \opt{--save} and \opt{--dry-run}. The exception
is a project file that names an output folder (chapter~\ref{ch:project}). All file names
and folder names on the command line start from the current folder.

\subsection*{What it prints}

\begin{itemize}
\item With an output tree, hdl-vet prints these lines, in this order:
  \begin{itemize}
  \item the line \code{wrote DIR}
  \item the line \code{saved FILE}, if you saved
  \item the line \code{N applied, M skipped}
  \item one line \code{changed PATH} for each changed file, with the name relative to the
    tree
  \item one line \code{added NAME} for each file that a script added
  \item the \code{note:} lines
  \end{itemize}
\item With \opt{--dry-run}, hdl-vet prints the diff on the standard output. It prints
  everything else on the error output.
\item If hdl-vet uses a project file, the first line, on the error output, is
  \code{project: PATH}.
\item If the run fails, hdl-vet prints the message on the error output.
\end{itemize}

\subsection*{Exit statuses}

This table is the one place that defines the exit statuses.

\begin{center}
\begin{tabular}{@{}lp{0.75\textwidth}@{}}
\toprule
Status & Meaning \\
\midrule
0 & The run worked. Without \opt{--dry-run}, hdl-vet writes the output tree. With \opt{--save}
  and no output, the run ends with 0 also when the sign-off failed so far
  (chapter~\ref{ch:session}). \\
1 & The sign-off failed (chapter~\ref{ch:signoff}). The report is on the error output, and
  hdl-vet writes no output tree. This holds with \opt{--dry-run} too. \\
2 & hdl-vet cannot do the request: a wrong command line, a file that it cannot read,
  a design that it cannot load at the start, a script that fails, a refusal of an
  operation, or drift. It writes nothing. \\
\bottomrule
\end{tabular}
\end{center}

Messages of the command line have two forms. The first form is for a command line that is
wrong in a way that the argument parser sees. Examples are an unknown option, a missing
value, and no command at all. In this case \cmd{vet} prints a usage line and then
\code{vet: error: ...}, and exits with status 2. For every other error, \cmd{vet} prints
one line, \code{vet: ...}, and exits with status 2.

\lstinputlisting[style=console]{examples/console/cli.txt}

In the transcript, look at these things:

\begin{itemize}
\item \cmd{vet --version} prints the program name and the version (here \code{<version>}).
\item The file list \file{flow.f} is outside the folder of the sources, so the
  command gives \opt{--base} \code{design}. The output tree has the file list at its top.
\item hdl-vet refuses the second run, because \file{out/flow} exists. The message names
  \opt{--force}. The third run uses it.
\item The transcript has five errors, each with exit status 2:
  \begin{itemize}
  \item no file list
  \item no output
  \item a file list that is not there
  \item a top module that is not in the design
  \item a script that cannot be read
  \end{itemize}
\end{itemize}

\subsection*{What replacing an output folder means}

\opt{--force} replaces the output folder. hdl-vet builds the new tree next to the
output folder. Then it renames the new tree to the name of the output folder. So a failed
run never leaves half a tree. This is the one place that describes the limits.

\begin{itemize}
\item With \opt{-o}, \opt{--force} replaces \emph{any} folder that you name,
  and this includes a folder that holds files of yours. Check the name.
\item It never replaces a folder inside the input tree, the input tree itself, or a
  folder that holds the input tree. It never replaces a folder that holds a script
  of the run or the project file.
\item If the output folder comes from the project file, \opt{--force} replaces it in
  two cases only:
  \begin{itemize}
  \item The folder is empty.
  \item The folder is the output of an earlier run and holds nothing else.
  \end{itemize}
  In the second case, \file{hdl-vet.json} must be a manifest of hdl-vet with a list of
  files. Every file in the folder must be that manifest or be in the list. Otherwise
  hdl-vet refuses. The message names a file that is not in the list.
  Output trees from before the list existed have no such list, so \opt{--force} refuses
  them.
\end{itemize}

\section{vet init}

\cmd{vet init DIR -f FILELIST} writes a starter project into \file{DIR}
(chapter~\ref{ch:first} and chapter~\ref{ch:project}).

\begin{description}
\item[\opt{-f}, \opt{--filelist} FILE] The file list of the design. It is required.
\item[\opt{--top} MODULE] The top module. If you give none, \cmd{vet init} takes the only
  top-level module of the design. If the design has several, \cmd{vet init} lists them
  (at most ten, the biggest hierarchy first) and asks for \opt{--top}.
\item[\opt{--base} DIR] The base of the file list.
\end{description}

\cmd{vet init} loads the design first, so it refuses a design that has errors. It
refuses a \file{DIR} that exists and is not empty. It refuses a \file{DIR} inside the
input tree, because the next run of the generator can delete it. When it refuses, it writes
nothing. It prints \code{created DIR}, exactly as you typed it, and one line
\code{wrote PATH} for each file.
